%% ARG Co-occurrence Network Under Climate Gradients: A Gene-Level Network Topology Study
%% Authors: James Angelo R. Dela Cruz, Angeleen Claire M. Lominoque, Hazelle N. Millarpis
%% Course: CMSC 191 NS -- Network Science
%%
%% Overleaf setup (copy-paste ready):
%%   1. Paste this entire file into main.tex
%%   2. Upload arg-bibliography.bib to the project root
%%   3. Create a folder named "figures" and upload:
%%        network_tropical.png, network_arid.png,
%%        network_temperate.png, centrality_by_zone.png

\documentclass[12pt,a4paper]{article}

\usepackage[margin=2.5cm]{geometry}%
\usepackage{authblk}%
\usepackage{graphicx}%
\usepackage{multirow}%
\usepackage{amsmath,amssymb,amsfonts}%
\usepackage{amsthm}%
\usepackage{booktabs}%
\usepackage{xcolor}%
\usepackage{hyperref}%
\usepackage[numbers,sort&compress]{natbib}%
\usepackage{setspace}%
\onehalfspacing%
\usepackage{parskip}%

\hypersetup{colorlinks=true,linkcolor=blue,citecolor=blue,urlcolor=blue}

\renewcommand\Affilfont{\small}

\begin{document}

\title{\textbf{Antibiotic Resistance Gene Co-occurrence Network Topology\\
Under Climate Gradients: A Gene-Level Network Analysis}}

\author[1]{James Angelo R. Dela Cruz}
\author[1]{Angeleen Claire M. Lominoque}
\author[1]{Hazelle N. Millarpis}
\affil[1]{Department of Physical Sciences and Mathematics,
  University of the Philippines Manila,
  670 Padre Faura St, Manila 1000, Philippines}

\date{}
\maketitle

%% ─── Abstract ─────────────────────────────────────────────────────────────────

\begin{abstract}
Antibiotic resistance genes (ARGs) in environmental metagenomes do not
accumulate in isolation; they are co-selected, co-transferred, and
co-regulated in ways that produce structured co-occurrence patterns.
Whether these patterns vary systematically with climate zone remains
largely unexplored.
This study constructs weighted ARG co-occurrence networks for three
major climate zones---tropical, arid, and temperate---using the
mARG dataset (Martiny \emph{et~al.}) comprising 119,232 metagenomic
samples, with climate zone assigned via the Beck \emph{et~al.}\ (2023)
K\"oppen-Geiger raster at $0.1^\circ$ resolution.
Pairwise Spearman correlations between 156 ARGs passing a 5\% presence
filter were used to define edges ($r \geq 0.3$, $p < 0.05$).
Louvain community detection revealed modularity scores of
$Q = 0.41$ (tropical), $Q = 0.26$ (arid), and $Q = 0.31$ (temperate),
indicating zone-specific variation in network compartmentalisation.
Hub ARGs identified by betweenness centrality differed across zones,
with \texttt{tet(M)} recurring in both tropical and arid networks.
A Kruskal-Wallis test on betweenness distributions across zones
yielded $H = 3.22$, $p = 0.20$ (not significant), suggesting that
while the \emph{identity} of hub ARGs differs by zone, the overall
structural centrality distribution is conserved.
These findings provide a network-level perspective on how climate context
may shape the ecology of antibiotic resistance.
\end{abstract}

\noindent\textbf{Keywords:} Antibiotic Resistance Genes, Co-occurrence Network,
K\"oppen-Geiger Climate Classification, Community Detection,
Louvain Algorithm, Betweenness Centrality, Network Topology, Environmental Metagenomics

%% ─── 1. Introduction ──────────────────────────────────────────────────────────

\section{Introduction}\label{intro}

The global rise of antibiotic resistance constitutes one of the most
urgent public-health challenges of the twenty-first century.
Antibiotic resistance genes (ARGs) are the molecular currency of this
crisis: they encode biochemical mechanisms---enzymatic inactivation,
efflux pumps, target modification---that render antibiotics
ineffective~\cite{WHO2023}.
Environmental metagenomics has demonstrated that ARGs are pervasive
across soils, sediments, and waterways, raising concern about the
``resistome'' as a reservoir from which clinically relevant resistance
elements may emerge~\cite{Allen2010}.

ARGs are rarely found in isolation.
They are frequently co-localised on mobile genetic elements such as
plasmids and integrons, creating co-selection pressure that drives
their co-occurrence in metagenomic samples~\cite{Partridge2018}.
Network science offers a natural framework for studying these
co-occurrence patterns.
By treating ARGs as nodes and statistically significant co-occurrences
as edges, one can characterise the global structure of resistance gene
ecology---identifying hub genes that link distinct resistance clusters,
and detecting communities of genes that tend to co-occur together.

Climate zone is a plausible predictor of ARG community structure.
Temperature and precipitation shape microbial community composition,
horizontal gene transfer rates, and the selective pressure from
antibiotic use in agriculture~\cite{vanBoeckel2014, Cavicchioli2019}.
Yet to our knowledge, no study has directly compared ARG co-occurrence
\emph{network topology} across climate zones using a large,
geographically diverse metagenomic dataset.

This study addresses that gap.
Using the mARG dataset---the largest unified ARG abundance table
compiled from published environmental metagenomic studies---we
assign climate zone labels via the high-resolution
K{\"o}ppen-Geiger raster of Beck \emph{et~al.}\ (2023), construct
weighted co-occurrence networks for tropical, arid, and temperate
zones, and compare their topology through community detection and
centrality analysis.

\section{Review of Related Literature}\label{rrl}

\subsection{The Global Burden of Antimicrobial Resistance}

Antimicrobial resistance has emerged as one of the defining public health
crises of the twenty-first century.
The World Health Organization classifies AMR among the top ten threats to
global health, and surveillance data indicate that resistant infections now
cause millions of deaths annually~\cite{WHO2023}.
Drug-resistant pathogens complicate the treatment of common bacterial
infections, surgical procedures, and cancer chemotherapy, eroding the
clinical foundations of modern medicine.
The burden is not distributed evenly: low- and middle-income countries,
which already face constrained healthcare infrastructure, carry a
disproportionate share of resistance-attributable mortality~\cite{LancetAMR2024}.

The environmental reservoir plays a central, if underappreciated, role in
this crisis.
Antibiotic resistance genes are ancient and pervasive features of
microbial genomes, predating the clinical antibiotic era by millions of
years~\cite{Allen2010}.
Soils, sediments, and water bodies harbour vast and diverse resistomes
whose composition reflects both natural selective pressures and the
accumulated legacy of antibiotic use in medicine, agriculture, and
aquaculture.
Environmental ARGs constitute a standing reservoir from which
clinically relevant resistance elements can be recruited into pathogens
through horizontal gene transfer~\cite{Partridge2018}.


\subsection{Climate Change as a Driver of Antimicrobial Resistance}

Growing evidence links climatic variables to the abundance, diversity,
and spread of environmental ARGs.
A systematic scoping review of the climate-AMR literature found consistent
associations between rising temperature and increased resistance rates across
multiple pathogen species and geographic
settings~\cite{SalgueiroCernuda2024}.
The biological mechanisms proposed to explain these associations include
accelerated bacterial growth at higher temperatures, increased rates of
horizontal gene transfer through bacterial conjugation, and heat-induced
stress responses that upregulate efflux pumps and other resistance
mechanisms~\cite{LancetWP2025}.

Extreme weather events compound these pressures in a different way.
Floods and heavy rainfall damage sanitation infrastructure and contaminate
water sources, releasing ARGs from agricultural and wastewater reservoirs
into the broader environment~\cite{LancetWP2025}.
Droughts concentrate resistant bacteria and their genetic cargo in
shrinking water bodies, increasing the probability of gene transfer events.
These pathways position climate change not merely as a background driver
of AMR but as an active amplifier of the spread and persistence of
resistance genes across environments and hosts~\cite{Cavicchioli2019}.

Despite this growing body of evidence, the relationship between climate
and the structural organisation of ARG communities at the gene-network
level remains largely uncharacterised.
Most existing studies examine ARG abundance or resistance rates as scalar
quantities, without capturing the co-occurrence relationships that reflect
the ecological and genetic linkages among resistance determinants.


\subsection{ARG Co-occurrence Networks in Environmental Metagenomics}

Network science has emerged as a productive framework for studying
the co-occurrence patterns of ARGs in environmental metagenomes.
By representing ARGs as nodes and statistically significant co-occurrence
events as edges, researchers can model the collective ecology of the
environmental resistome, rather than examining individual genes in
isolation.
Early applications of this approach in diverse environments, including
soil, freshwater, and marine sediments, established that ARG co-occurrence
networks share structural properties with other biological networks:
sparse but non-trivial connectivity, modular organisation into gene
communities, and the presence of highly connected hub genes that link
otherwise distinct clusters~\cite{ZhuChen2013}.

The co-occurrence of ARGs in the same environmental sample most commonly
reflects physical co-localisation on shared mobile genetic elements such
as plasmids, integrons, and transposons~\cite{Partridge2018}.
When two genes are repeatedly found together across geographically diverse
samples, it suggests that they are co-selected and co-transferred as a
functional unit.
This makes co-occurrence network topology an indirect but tractable proxy
for the genetic architecture of the resistome.
Subsequent studies applying network analysis to global metagenomic datasets
confirmed that co-occurrence patterns are highly context-dependent:
ARG communities in sewage, agricultural soil, and pristine environments
show distinct structural profiles, consistent with the influence of
antibiotic selection pressure and microbial community composition on
network organisation~\cite{Martiny2022}.

The use of large, unified ARG abundance tables, such as the mARG dataset
compiled by Martiny \emph{et~al.}, makes it possible to extend this
approach to the climate gradient question.
Analysing co-occurrence structure across samples grouped by climate zone
allows the structural properties of the resistome network to be compared
under systematically different environmental conditions, complementing
and extending the scalar abundance analyses that dominate the current
climate-AMR literature.


\subsection{Gap in the Literature}

Despite parallel advances in climate-AMR research and resistome network
analysis, no published study has directly compared ARG co-occurrence
network topology across major climate zones using a large, geographically
diverse, and uniformly processed metagenomic dataset.
Existing network studies are typically restricted to single environment
types or single geographic regions.
Existing climate-AMR studies focus on resistance rates rather than
network structure.
The present study addresses this intersection, asking not simply whether
climate zone correlates with ARG abundance, but whether it reshapes the
topological organisation of the resistance gene co-occurrence network.


%% ─── 2. Methods ───────────────────────────────────────────────────────────────

\section{Methods}\label{methods}

\subsection{Dataset}\label{sec:data}

We used the mARG dataset published by Martiny \emph{et~al.}\ (Zenodo
accession 6919377), which aggregates ARG abundance tables from 119,232
environmental metagenomic samples.
The dataset provides, for each sample-ARG pair, the fragment count
aligned to a reference ARG sequence (\texttt{fragmentCountAln}),
the reference sequence length, and the total number of trimmed reads
(\texttt{trimmed\_fragments}).
Geographic coordinates were parsed from the accompanying
\texttt{metadata.tsv} file, which records sample location in
degree-minute format (e.g.\ ``18.05 S 47.03 E'').

Abundance was normalised to Reads Per Kilobase per Million mapped reads
(RPKM):
\begin{equation}
  \text{RPKM}_{ij} =
    \frac{\text{fragmentCountAln}_{ij} \times 10^{6}}
         {\left(\text{refLength}_{j} / 1000\right)
          \times \text{trimmedFragments}_{i}}
  \label{eq:rpkm}
\end{equation}
where $i$ indexes samples and $j$ indexes ARG sequences.
This normalisation accounts for both gene length bias and variation in
sequencing depth.
The resulting wide-format matrix has dimensions
$119{,}232 \text{ samples} \times 2{,}999 \text{ ARGs}$.

\subsection{Climate Zone Assignment}\label{sec:zones}

Each sample was assigned to a major climate class using the
Beck \emph{et~al.}\ (2023) K{\"o}ppen-Geiger raster at
$0.1^\circ$ spatial resolution (Figshare dataset 21789074).
The raster encodes 30 climate types as integer codes; we mapped
these to three analytical zones:
\begin{itemize}
  \item \textbf{Tropical} (codes 1--3, class A)
  \item \textbf{Arid} (codes 4--7, class B)
  \item \textbf{Temperate} (codes 8--28, classes C and D merged)
  \item Polar (codes 29--30, class E) --- excluded
\end{itemize}
Point queries were performed in Python using \texttt{rasterio}.
Samples with coordinates outside the raster extent or returning a
no-data value were discarded.
This yielded 52,064 samples across the three zones
(Table~\ref{tab:zones}).

\begin{table}[htbp]
  \centering
  \caption{Sample counts per climate zone after K{\"o}ppen-Geiger
           assignment.}
  \label{tab:zones}
  \begin{tabular}{lcc}
    \toprule
    Climate Zone & Samples & Fraction (\%) \\
    \midrule
    Temperate    & 44,601  & 85.7 \\
    Tropical     &  4,443  &  8.5 \\
    Arid         &  3,020  &  5.8 \\
    \midrule
    \textbf{Total} & \textbf{52,064} & \textbf{100.0} \\
    \bottomrule
  \end{tabular}
\end{table}

\subsection{ARG Filtering}\label{sec:filter}

ARGs present in fewer than 5\% of samples (i.e.\ fewer than
$0.05 \times 52{,}064 \approx 2{,}603$ samples) were removed as
insufficiently prevalent to construct meaningful co-occurrence edges.
This reduced the feature set from 2,999 to \textbf{156 ARGs}.

\subsection{Co-occurrence Network Construction}\label{sec:network}

For each climate zone, we computed all pairwise Spearman rank
correlations between the 156 retained ARG abundance vectors.
An undirected, weighted edge was added between ARG $i$ and ARG $j$
if and only if:
\begin{equation}
  r_s(i, j) \geq 0.3 \quad \text{and} \quad p < 0.05
  \label{eq:edge}
\end{equation}
where $r_s$ is the Spearman correlation coefficient and $p$ is its
two-tailed $p$-value.
The threshold $r \geq 0.3$ follows established practice in
microbiome co-occurrence network studies~\cite{Faust2012} and
excludes weakly correlated pairs that may reflect chance associations.
Only positive correlations were retained, as negative correlations
are ecologically ambiguous in this context (competitive exclusion
\emph{vs.} differential habitat preference).
Networks were implemented as weighted undirected graphs in Python
using \texttt{NetworkX} and exported as \texttt{.gexf} files.

\subsection{Community Detection}\label{sec:community}

Community structure was assessed using two algorithms:
\begin{enumerate}
  \item \textbf{Louvain algorithm} (primary)~\cite{Blondel2008}:
    a greedy modularity-maximisation heuristic that iteratively
    reassigns nodes to communities to improve the modularity score
    $Q$, defined as
    \begin{equation}
      Q = \frac{1}{2m} \sum_{ij}
          \left[ A_{ij} - \frac{k_i k_j}{2m} \right]
          \delta(c_i, c_j)
      \label{eq:modularity}
    \end{equation}
    where $A_{ij}$ is the edge weight, $k_i$ is the weighted degree
    of node $i$, $m$ is the total edge weight, $c_i$ is the community
    assignment of node $i$, and $\delta$ is the Kronecker delta.
    Values of $Q > 0.3$ are conventionally regarded as indicating
    meaningful community structure~\cite{Newman2006}.
  \item \textbf{Girvan-Newman algorithm} (comparison)~\cite{Girvan2002}:
    a divisive hierarchical approach that removes edges with the
    highest betweenness centrality iteratively; the first-level
    partition was used.
\end{enumerate}

\subsection{Centrality Metrics}\label{sec:centrality}

Three node-level centrality measures were computed per zone:
\begin{enumerate}
  \item \textbf{Degree centrality}: fraction of other nodes to which
    a given ARG is connected.
  \item \textbf{Betweenness centrality}: proportion of shortest paths
    in the network that pass through a given ARG, normalised by the
    total number of pairs.
    Hub ARGs with high betweenness act as bridges between co-occurrence
    communities.
  \item \textbf{Clustering coefficient}: probability that two
    neighbours of a given ARG are themselves connected,
    reflecting local clustering.
\end{enumerate}

\subsection{Cross-zone Comparison}\label{sec:stats}

To test whether betweenness centrality distributions differ
significantly across climate zones, we applied the
Kruskal-Wallis $H$-test~\cite{Kruskal1952}, a non-parametric
one-way analysis of variance appropriate for non-normal distributions.
Significance was set at $\alpha = 0.05$.

%% ─── 3. Results ───────────────────────────────────────────────────────────────

\section{Results}\label{results}

\subsection{Network Topology}\label{sec:res-topology}

All three climate zone networks contained 155 nodes (one ARG with
no significant correlations was isolated in each zone and contributed
no edges).
Network statistics are summarised in Table~\ref{tab:netstat}.

\begin{table}[htbp]
  \centering
  \caption{Summary topology of ARG co-occurrence networks per climate zone.}
  \label{tab:netstat}
  \begin{tabular}{lcccc}
    \toprule
    Zone      & Edges & Density & Avg.\ Clustering & Components \\
    \midrule
    Tropical  & 1,730 & 0.145   & 0.315            & 24 \\
    Arid      & 2,292 & 0.192   & 0.326            & 15 \\
    Temperate & 1,960 & 0.164   & 0.323            & 18 \\
    \bottomrule
  \end{tabular}
\end{table}

The arid zone network was the most densely connected (density $= 0.192$)
with the fewest connected components (15), suggesting a more integrated
resistance gene ecology.
The tropical network was the sparsest (density $= 0.145$) with the most
components (24), consistent with a more partitioned co-occurrence
structure.
Network visualisations are shown in Figure~\ref{fig:networks}.

\begin{figure}[htbp]
  \centering
  \includegraphics[width=0.92\linewidth]{network_tropical.png}
  \caption{(a) \textbf{Tropical Climate Zone} --- 1,730 edges, $Q = 0.41$,
    Louvain communities: 27.
    Nodes are coloured by community and sized by betweenness centrality;
    labels shown for the top-12 hub ARGs by betweenness.
    Edge width is proportional to Spearman $r$.}
  \label{fig:network_tropical}
\end{figure}

\begin{figure}[htbp]
  \centering
  \includegraphics[width=0.92\linewidth]{network_arid.png}
  \caption{(b) \textbf{Arid Climate Zone} --- 2,292 edges, $Q = 0.26$,
    Louvain communities: 20.}
  \label{fig:network_arid}
\end{figure}

\begin{figure}[htbp]
  \centering
  \includegraphics[width=0.92\linewidth]{network_temperate.png}
  \caption{(c) \textbf{Temperate Climate Zone} --- 1,960 edges, $Q = 0.31$,
    Louvain communities: 24.}
  \label{fig:networks}
\end{figure}


\subsection{Community Detection}\label{sec:res-community}

Louvain community detection results are presented in
Table~\ref{tab:modular}.
The tropical zone exhibited the highest modularity ($Q = 0.41$),
indicating strongly demarcated ARG communities.
The arid zone fell below the conventional threshold of 0.3
($Q = 0.26$), suggesting a more integrated, less compartmentalised
co-occurrence network.
The temperate zone was intermediate ($Q = 0.31$).
Girvan-Newman partitioning produced qualitatively similar community
counts.

\begin{table}[htbp]
  \centering
  \caption{Community detection results per climate zone.}
  \label{tab:modular}
  \begin{tabular}{lccc}
    \toprule
    Zone      & Louvain Communities & Modularity $Q$ & GN Communities \\
    \midrule
    Tropical  & 27 & \textbf{0.4135} & 25 \\
    Arid      & 20 & 0.2601          & 16 \\
    Temperate & 24 & 0.3083          & 19 \\
    \bottomrule
  \end{tabular}
\end{table}

\subsection{Hub ARGs by Betweenness Centrality}\label{sec:res-hubs}

Table~\ref{tab:hubs} lists the top five ARGs by betweenness centrality
in each zone.
\texttt{tet(M)} variants appear as primary hubs in both the tropical
and arid zones, suggesting that this mobile tetracycline-resistance
gene occupies a structurally central bridging role across ecologically
distinct environments.
In contrast, the temperate zone is dominated by \texttt{lnu(B)} and
\texttt{aph(6)-Id}, genes conferring lincosamide and aminoglycoside
resistance respectively, which were not among the top hubs in tropical
or arid networks.
\texttt{lsa(E)}, conferring streptogramin and lincosamide resistance,
emerged as a top hub in both arid and temperate zones.

\begin{table}[htbp]
  \centering
  \caption{Top 5 hub ARGs by betweenness centrality per climate zone.}
  \label{tab:hubs}
  \begin{tabular}{llcc}
    \toprule
    Zone      & ARG                      & Betweenness & Degree Centr. \\
    \midrule
    \multirow{5}{*}{Tropical}
      & \texttt{tet(M)\_13\_AM990992} & 0.060 & 0.299 \\
      & \texttt{tet(Q)\_4\_Z21523}    & 0.045 & 0.325 \\
      & \texttt{tet(M)\_12\_FR671418} & 0.044 & 0.247 \\
      & \texttt{erm(B)\_10\_U86375}   & 0.039 & 0.143 \\
      & \texttt{tet(O)\_3\_Y07780}    & 0.030 & 0.344 \\
    \midrule
    \multirow{5}{*}{Arid}
      & \texttt{aadA2\_1\_NC\_010870}     & 0.081 & 0.182 \\
      & \texttt{lsa(E)\_1\_JX560992}      & 0.063 & 0.240 \\
      & \texttt{tet(M)\_13\_AM990992}     & 0.061 & 0.331 \\
      & \texttt{lnu(A)\_1\_M14039}        & 0.056 & 0.032 \\
      & \texttt{aph(3'')-Ib\_5\_AF321551} & 0.040 & 0.182 \\
    \midrule
    \multirow{5}{*}{Temperate}
      & \texttt{lnu(B)\_2\_JQ861959}   & 0.094 & 0.273 \\
      & \texttt{aph(6)-Id\_1\_M28829}  & 0.079 & 0.149 \\
      & \texttt{aph(3')-Ia\_7\_X62115} & 0.054 & 0.071 \\
      & \texttt{mef(A)\_2\_U83667}     & 0.050 & 0.214 \\
      & \texttt{lsa(E)\_1\_JX560992}   & 0.048 & 0.305 \\
    \bottomrule
  \end{tabular}
\end{table}

\subsection{Cross-zone Centrality Comparison}\label{sec:res-kw}

The Kruskal-Wallis test on betweenness centrality distributions
across the three zones yielded $H = 3.219$, $p = 0.200$
(Figure~\ref{fig:boxplot}).
This result is not statistically significant at $\alpha = 0.05$,
indicating that the \emph{distribution} of betweenness centrality
values is statistically similar across zones.
The apparent zone differences in hub ARG identity (Section~\ref{sec:res-hubs})
therefore reflect shifts in \emph{which} ARGs occupy central positions,
rather than a change in the overall degree of network centralisation.

\begin{figure}[htbp]
  \centering
  \includegraphics[width=0.85\linewidth]{figures/centrality_by_zone}
  \caption{Distribution of betweenness centrality across climate zones.
    Kruskal-Wallis test: $H = 3.22$, $p = 0.20$ (not significant).
    The three zones share a similar right-skewed centrality distribution;
    most ARGs have low betweenness while a small number serve as hubs.}
  \label{fig:boxplot}
\end{figure}

%% ─── 4. Discussion ────────────────────────────────────────────────────────────

\section{Discussion}\label{discussion}

\subsection{Climate-stratified ARG network structure}

All three co-occurrence networks exhibit properties consistent with
other biological co-occurrence networks: sparse but non-trivial density,
high clustering coefficients, and modular organisation.
The zone-level differences in modularity are particularly informative.

The tropical network's high modularity ($Q = 0.41$) is consistent with
the ecological complexity of tropical environments.
Tropical climates support high microbial biodiversity, complex soil
chemistry, and a wide range of ecological niches spanning rainforest
soils, freshwater systems, and coastal sediments.
This ecological heterogeneity may produce compartmentalised ARG
communities because different microhabitats within tropical zones impose
distinct selective pressures, favouring different sets of resistance
determinants and mobile genetic elements~\cite{Cavicchioli2019}.
When samples from ecologically distinct niches are pooled within a single
climate zone, the resulting network naturally separates into communities
corresponding to those niche-specific resistomes.
The high modularity value observed here is thus broadly consistent with
tropical zones harbouring a more partitioned, ecologically structured
resistome.

The comparatively low modularity of the arid zone ($Q = 0.26$, below
the conventional 0.3 threshold) points toward a different pattern of
ARG ecology.
Arid environments impose a relatively uniform set of stressors:
water limitation, high temperatures, and low organic carbon.
These conditions may select for a more generalist microbial community
with a correspondingly less compartmentalised set of resistance
mechanisms~\cite{SalgueiroCernuda2024}.
An alternative explanation is that arid-zone samples in the mARG
dataset are dominated by a small number of environment types, such as
desert soil and dryland agricultural sites, which would reduce
between-niche variation and compress the network into fewer, less
distinct communities.
Either way, the low modularity of the arid network suggests that
resistance genes in dry climates co-occur in a more integrated,
less segregated fashion.

The temperate zone occupies an intermediate position ($Q = 0.31$),
reflecting the environmental diversity of the C and D climate classes,
which span Mediterranean shrublands, continental grasslands, and
boreal forest soils across multiple continents.
The temperate zone is also heavily influenced by agricultural antibiotic
use, particularly in North America and Europe, where animal husbandry
contributes substantially to the resistance gene load in soils and
waterways~\cite{vanBoeckel2014}.
This anthropogenic input may reinforce certain co-occurrence edges while
diluting the natural ecological signal, yielding the intermediate
modularity observed.


\subsection{Hub ARG identity across zones}

The recurrence of \texttt{tet(M)} as a top hub in both tropical and
arid zones reflects its well-documented ecological breadth.
\texttt{tet(M)} encodes a ribosomal protection protein that confers
resistance to tetracyclines, one of the most widely used antibiotic
classes in both human medicine and animal production worldwide.
Critically, \texttt{tet(M)} is carried on conjugative transposons of
the Tn916 family, which are capable of transfer across a wide range of
Gram-positive and Gram-negative bacterial genera~\cite{Roberts2011}.
This broad-host-range mobility gives \texttt{tet(M)} the ecological
capacity to co-occur with a diverse set of other resistance genes across
phylogenetically distant bacterial hosts, explaining its structural
centrality in networks representing ecologically distinct environments.
Its recurrence across both tropical and arid zones, despite the
substantially different community compositions and modularity structures
of those networks, underscores that certain highly mobile ARGs can
function as network connectors irrespective of local climate conditions.

In contrast, the temperate zone hub profile is dominated by
\texttt{lnu(B)} and \texttt{aph(6)-Id}, genes conferring resistance
to lincosamides and aminoglycosides respectively.
Neither gene appeared among the top five hubs in the tropical or arid
networks.
Lincosamide and aminoglycoside antibiotics are both extensively used in
clinical and veterinary settings concentrated in temperate-zone countries,
particularly across Europe and North America~\cite{vanBoeckel2014}.
The emergence of these genes as network hubs in temperate samples is
broadly consistent with selection pressure from heavy antibiotic use
reinforcing co-selection and co-occurrence linkages among resistance
determinants adapted to high-usage environments.
This pattern suggests that the structural centralisation of specific
ARGs in co-occurrence networks may in part reflect the regional
antibiotic usage profile of the dominant sampling locations within
a climate zone.

\texttt{lsa(E)}, a gene conferring combined resistance to
pleuromutilins, lincosamides, and streptogramins, originally
characterised in \emph{Enterococcus} isolates~\cite{Lozano2012},
appeared as a top hub in both arid and temperate zones.
This cross-zone prevalence is notable because the two zones differ
substantially in network density, modularity, and microbial community
context.
One plausible explanation is that \texttt{lsa(E)} resides on
broad-host-range plasmids that facilitate its dissemination across
phylogenetically and ecologically diverse microbial communities,
making it a persistent co-occurrence partner for other resistance genes
regardless of climate context.
The clinical relevance of \texttt{lsa(E)} is increasing, as it has been
detected in livestock-associated \emph{Staphylococcus} and
\emph{Enterococcus} strains in multiple geographic regions, and its
identification as a cross-zone hub in environmental samples warrants
attention in AMR surveillance frameworks.

Taken together, the zone-specific hub profiles suggest two
coexisting processes in the environmental resistome.
First, highly mobile genes with broad host ranges (exemplified by
\texttt{tet(M)}) generate stable co-occurrence hubs that are
structurally central across climatically diverse networks.
Second, the identity of additional hubs is shaped by local conditions,
including antibiotic use patterns and ecological niche diversity,
producing zone-specific centrality signatures.
This distinction has practical implications: surveillance strategies
targeting structurally central ARGs should prioritise broadly mobile
genes as universal targets while also monitoring for the locally
dominant hubs that characterise each climate or geographic context.

\subsection{Interpretation of non-significant Kruskal-Wallis result}

The non-significant Kruskal-Wallis test ($p = 0.20$) should not be
interpreted as evidence that climate zone has no effect on ARG
network structure.
Network-level comparison via overall centrality distribution is a
coarse measure; it cannot capture the zone-specific shifts in hub ARG
identity documented in Table~\ref{tab:hubs}.
Moreover, the heavy overrepresentation of temperate samples
(85.7\% of the zoned dataset) may have suppressed inter-zone variation
by biasing the temperate distribution with very high statistical power
relative to tropical and arid.
Future work with more balanced zone sampling would provide a more
definitive test.

\subsection{Limitations}

Several limitations bear noting.
First, the dataset is heavily skewed toward temperate-zone samples;
tropical and arid zones each contribute fewer than 10\% of the zoned
data, which limits the statistical power of inter-zone comparisons.
Second, the 5\% presence filter, while necessary for network
stability, excluded 94.8\% of ARGs, potentially omitting
zone-specific rare resistance elements.
Third, the Spearman correlation threshold ($r \geq 0.3$) and
$p$-value cutoff ($< 0.05$) are conventional but arbitrary; network
topology is sensitive to these choices.
Finally, co-occurrence does not imply physical co-localisation on
a mobile genetic element; confounding by shared ecological niche
cannot be excluded without confirmatory genomic evidence.

%% ─── 5. Conclusion ────────────────────────────────────────────────────────────

\section{Conclusion}\label{conclusion}

This study presents the first large-scale, climate-stratified
co-occurrence network analysis of environmental ARGs.
Using the mARG dataset and K{\"o}ppen-Geiger climate zone assignment,
we constructed weighted ARG co-occurrence networks for tropical, arid,
and temperate zones and characterised their topology through community
detection and centrality analysis.

Key findings are: (1) all three zones exhibit non-trivial modular
structure; (2) tropical networks are the most compartmentalised
($Q = 0.41$), while arid networks are the most integrated
($Q = 0.26$); (3) \texttt{tet(M)} and \texttt{lsa(E)} recur as
structural hubs across multiple zones, consistent with their known
mobility; and (4) overall betweenness centrality distributions are
statistically similar across zones despite qualitative differences
in hub ARG identity.

These findings contribute to our understanding of how climate context
shapes the ecology of antibiotic resistance at the gene-network level
and provide a framework for future comparative resistome studies with
more geographically balanced sampling.

From a public health and surveillance standpoint, the identification of
\texttt{tet(M)} and \texttt{lsa(E)} as cross-zone structural hubs
suggests that monitoring these genes in environmental samples may
provide early warning signals of resistome shifts across climatically
diverse settings.
The zone-specific hub profiles further indicate that universal
surveillance frameworks should be complemented by regionally tailored
monitoring that accounts for the local antibiotic use patterns and
ecological conditions that shape which resistance genes occupy central
positions in the co-occurrence network.
Future work with geographically balanced sampling and longitudinal
datasets will be needed to determine whether the structural differences
observed here are stable features of climate-zone resistomes or reflect
sampling artefacts in the current mARG dataset.
%% ─── Back matter ──────────────────────────────────────────────────────────────

\section*{Data and Code Availability}

The mARG dataset used in this study is publicly available at Zenodo
(accession 6919377).
The K\"oppen-Geiger climate raster is available at Figshare
(dataset 21789074; Beck \emph{et~al.}\ 2023).
All analysis scripts are available in the project repository at
\url{https://github.com/Jamesardc19/pitch2-arm-gene-network-main}.

\section*{Author Contributions}

J.A.R.D.C.\ designed and implemented the full analytical pipeline,
including data preprocessing, network construction, community
detection, centrality analysis, and visualisation.
A.C.M.L.\ contributed to dataset sourcing and literature review for
the Background and Methods sections.
H.N.M.\ conducted the literature review and drafted the Introduction
and Review of Related Literature.
All authors contributed to the Conclusion.

%% ─── Bibliography ─────────────────────────────────────────────────────────────

\newpage
\bibliographystyle{unsrt}%
\bibliography{arg-bibliography}%

\end{document}
